\documentclass[a4paper,10pt]{article}
\usepackage[utf8]{inputenc}
\usepackage[T1]{fontenc}
\usepackage[french]{babel}
\usepackage{amsmath,amssymb}
\usepackage{geometry}
\usepackage{array,tabularx}
\usepackage{tikz}
\geometry{margin=1.6cm}
\pagestyle{empty}
\newcommand{\ligne}{\par\vspace{0.55cm}\noindent\dotfill\par}
\newcommand{\carreaux}[1]{\begin{center}\begin{tikzpicture}\draw[step=0.5cm,gray!35,very thin] (0,0) grid (17,#1);\end{tikzpicture}\end{center}}
\newenvironment{exercice}[2]{\paragraph{Exercice #1 \hfill #2 points}\mbox{}\par}{\medskip}
\begin{document}
\begin{center}\Large\bfseries Devoir surveillé 18 - Les probabilités\end{center}
\vspace{2mm}
\begin{exercice}{1}{4}On tire une carte dans un jeu de $32$ cartes. Calculer la probabilité d'obtenir un roi, puis un coeur.\end{exercice}\begin{exercice}{2}{5}Une roue comporte $8$ secteurs égaux : $3$ rouges, $2$ bleus et $3$ verts. Calculer les probabilités des trois couleurs.\end{exercice}\begin{exercice}{3}{5}Une urne contient $4$ boules noires, $5$ blanches et $1$ rouge. Quelle est la probabilité de ne pas tirer une boule rouge ?\end{exercice}\begin{exercice}{4}{6}Inventer une expérience dont une issue a pour probabilité $\frac{3}{10}$.\end{exercice}\end{document}
